\documentclass{article}  % good style
\usepackage[a4paper]{geometry}  % wider margin
\usepackage{graphicx}  % for images
\usepackage[T1]{fontenc}
\usepackage{float}  % used for image positioning, very handy.
\usepackage{gensymb} % has the degree symbol
\usepackage{amsmath} % math
\usepackage{listings} % code
\usepackage{hyperref} % hyperlinks; \href{link_here}{label_here} -> opposite to md

\begin{document}
	
	\begin{center}
		\begin{Large}
			Brendan McCann | Masters Thesis Two-Pager
		\end{Large}
	\end{center}
	
    \subsection*{Research question \& subsequent hypothesis}

    The provisional title for my thesis is: \textbf{\emph{What Course Should I study?} Investigating the application of a Recommender System to assist in college course counselling.}
    The firstresearch question to investigate is: \emph{How can computational methods (e.g. AI) be used to assist in the Guidance Counselling (GC) process?}
    My hypothesis is that a Recommender System (RS) can be leveraged to suggest strong-fit college courses to students, 
    which would help their college course research.
    It is very important to note that this RS would not replace the role of GC's, nor will it replace the college course research process that every student must undertake.
    Rather, the RS would \textbf{augment} the college course research process.
    Also, the recommendations should \textbf{not be seen as conclusive or final} - they should simply be seen as \textbf{suggestions} that act as 
    \textbf{starting points} for the student's college course research process.

    \subsection*{Methodology \& current challenges}

    One challenge I am currently investigating is:
    \emph{What data do we gather from the student in order to make informed and effective course recommendations?}
    Both GC's I talked to mentioned that a lot of information can be learned about the student when you ask what LC subjects they like/dislike and why.
    This seemed promising to me, and I will likely integrate this into our RS.

    I am considering integrating a vocational interest quiz, likely the \href{https://openpsychometrics.org/tests/RIASEC/}{Holland Codes Quiz}. 
    However, one of the GC's I talked to expressed concerns that Holland codes provide narrow snapshots of the students and can be very limiting.
    I do believe her concerns are valid, but the goal of this RS is not to provide perfect-fit college courses, 
    it is to provide starting points for the student.
    It is also worth mentioning that John Holland's theories are ubiquitous in GC.
    So, I would still lean towards utilising Holland Codes in my RS.
    
    The second challenge I am currently investigating is: \emph{Do I utilise a traditional-style RS with collaborative/content-based filtering, or 
    an LLM-based RS?} The advantage of a traditional-style RS lies in its simplicity - it will be easy to decipher why a certain recommendation was made,
    which is important for scrutability, model transparency and user trust. 
    Although, this simplicity is also its biggest disadvantage - the questions asked to the student to gather data would have to be close-ended (\emph{e.g. Do you like Physics, yes or no?}), 
    which can be at risk for over-simplifying the GC process.

    With LLM's, the data-gathering process can be in the style of open-ended questions, for example: \emph{Tell me why you like or dislike Physics?}
    The ability of LLM's to understand natural language lends itself well to the GC process, because GC is a very human-centered process, 
    often in the form of conversations. However, my main concern with an LLM-based RS is a lack of control. 
    These are incredibly complex systems that are black-boxes by nature, and I don't like the idea of off-handing all the algorithmic complexity
    to an LLM. 
    To be really frank, it feels a bit lazy, and I would like to be more ambitious in my implementation.

    Perhaps a hybrid RS can be used - after all, that is what Netflix is doing today. 
    A non-LLM-based RS can trim down the list of potential recommendations to a smaller subset (called candidate generation),
    and then this smaller subset is fed to an LLM-based RS which can make sophisticated recommendations beyond the capabilities of 
    a traditional RS. This is something I would like to investigate further.

    \newpage

    \subsection*{Evaluation}

    The one audience I wish to evaluate my RS with is definitely students. 
    As mentioned previously, I would like my RS to provide starting points to the student, not conclusive recommendations.
    With this in mind, I would like my RS to give \textbf{diverse} recs, ensuring a broad range of course areas suggested.
    The RS should also aid the students course research process by providing them with courses they may not have thought about 
    otherwise. In other words, the RS should provide the student with \textbf{novel} and \textbf{serendipitous} recs.
    Diversity, novelty and serendipity are well-known features of an effective RS, and there are established methods of evaluating these 
    features in an RS. For example, they are outlined in Chapter 7 of Recommender Systems - The Textbook by Charu Aggarwal.

    Another audience I wish to evaluate my RS with is Guidance Counsellors. 
    I have not given a lot of thought as to how I would evaluate my RS with GC's, but off the top of my head,
    we could provide the GC with a sample student's Holland Code scores (RIASEC) and their liked/disliked LC subjects. 
    We can then ask them to provide a list of courses they would suggest to this student's profile, and then ask the student themself
    to compare the GC's results with my RS' suggestions.

    In terms of benchmarks, I believe a good baseline is to feed the information gathered about the student into an LLM model, such as ChatGPT,
    and then the student can compare the recommendations of:

    \begin{enumerate}
        \item A real GC
        \item My RS
        \item An LLM like ChatGPT
    \end{enumerate}
	
\end{document}
